
\chapter{FORMA DE PREZENTARE.\\ STRUCTURA LUCRĂRII}

Materialele se prezintă editate cu toate semnele diacritice din limba română, cu figurile și tabelele incluse în text. Se va folosi fontul Times New Roman, mărime 12 pt. Formatul paginii va fi A4; margini: 2.5 cm – sus, jos, stânga, dreapta.

Titlurile capitolelor se vor scrie centrat, cu majuscule și vor fi precedate de un simbol numeric, începând cu cifra 1, urmată de un punct; pentru subtitluri, vor urma și alte cifre, despărțite între ele prin puncte, în funcție de gradul titlului subordonat (ex.: 1. CAPITOL; 1.1. Subcapitol; 1.1.1. Subcapitol secundar). Nu se vor utiliza mai mult de 3 grade de titluri (capitole, subcapitole, subcapitole secundare).

Fonturi și aliniere pentru titluri:
\begin{itemize}
	\item titlul capitolului: font 18 pt, \textbf{bold}, majuscule, centrat;
	\item titluri subcapitole: font 16 pt, \textbf{bold}, \textit{Sentence case}, aliniat la stânga, tab de 1 cm;
	\item titluri subcapitole secundare: font 14 pt, \textbf{bold}, \textit{Sentence case}, aliniat la stânga, tab de 1 cm.
\end{itemize}

 Pentru spațierea titlurilor, se vor lăsa rânduri libere (font 12 pt) astfel:
 \begin{itemize}
	\item pentru titlul capitolului (pe prima pagină): 5 rânduri deasupra, 3 rânduri dedesubt;
	\item pentru subtitluri: 2 rânduri deasupra, 1 rând dedesubt;
	\item pentru subtitlurile secundare: 1 rând deasupra, 1 rând dedesubt.
\end{itemize}

După titluri (capitole, subcapitole etc.) nu se pune punct; excepție fac titlurile urmate de subtitluri (ex.: „Introducere. Aspecte fundamentale în ...”).

Un subtitlu trebuie să fie urmat întotdeauna de cel puțin 2 rânduri de text pe aceeași pagină; dacă nu a rămas suficient spațiu, se va deplasa totul pe pagina următoare. \\

Textul propriu-zis se va scrie cu font Times New Roman, spațiere 1.15\textbf{--}1.2 rânduri, mărime 12 pt, \textit{Justified}, tab de 1 cm la primul rând al paragrafului.

În text, la sfârșit de rând, dacă este cazul, se face despărțirea cuvintelor în silabe, pentru a nu lăsa spații albe, neuniforme între cuvinte.

Se recomandă să nu se folosească pe parcursul lucrării multe tipuri de simboluri pentru enumerări (\textit{Bullets}); enumerările se vor scrie cu un tab de 1.5 cm, apoi cele imbricate cu tab de 2.25 cm, respectiv 3 cm. De asemenea, se recomandă folosirea cratimei (-) pentru cuvintele compuse și a liniei lungi (\textbf{--}, \textit{En dash}) pentru enumerări sau explicații.

Se recomandă să nu existe în partea de jos a paginii un paragraf nou cu mai puțin de 2 rânduri; același lucru pentru un paragraf care se termină la începutul paginii următoare.

Structura lucrării trebuie să cuprindă:
\begin{itemize}
    \item pagini inițiale (înaintea corpului lucrării):
    \begin{itemize}
        \item pagină de titlu (pagină goală pe verso) \textbf{--} se vor înlocui secțiunile marcate cu galben din prezentul ghid; dacă titlul ocupă 1 rând, se vor lăsa libere rândurile de deasupra și de dedesubt; dacă titlul ocupă 2 rânduri, se va lăsa liber rândul de dedesubt; titlurile coordonatorilor științifici pot fi Prof. dr. ing., Conf. dr. ing., Ș.L. dr. ing. sau Dr. ing. / Ing. (externi UPB);
	\item Anexa 2 \textbf{--} se va descărca de pe site-ul de absolvire și se va include în continuare, conform specificației prezentului ghid;
	\item declarație de onestitate academică (pagină goală pe verso) – se vor înlocui secțiunile marcate cu galben din prezentul ghid;
	\item cuprins \textbf{--} se vor înlocui secțiunile marcate cu galben din prezentul ghid;
	\item lista tabelelor, lista figurilor și lista abrevierilor (în această ordine) \textbf{--} se vor înlocui secțiunile marcate cu galben din prezentul ghid;
    \end{itemize}
    \item corpul lucrării (capitole numerotate):
	\begin{itemize}
	    \item introducere (capitolul 1) \textbf{--} domeniul abordat, descrierea problemei, evidențierea oportunităților de îmbunătățire, structura concretă a lucrării;
	\item următoarele capitole propriu-zise ale lucrării (capitolele 2, 3, ...) \textbf{--} acestea pot fi, spre exemplu, un capitol de trecere în revistă a stării artei și de detaliere concisă a noțiunilor teoretice (tehnici de analiză și prelucrare de semnale, modele de învățare automată etc.) utilizate în cadrul lucrării (capitolul 2), unul sau două capitole de definire și detaliere a unui sistem sau a unor tehnici originale propuse (capitolul 3 / capitolele 3-4) și un capitol de prezentare a seturilor de date utilizate, a metodologiei experimentale și a rezultatelor obținute, împreună cu interpretarea acestora (capitolul 4/5);
	\item concluzii (ultimul capitol) \textbf{--} concluzii generale, contribuții originale, dezvoltări ulterioare;
    \end{itemize}
    \item bibliografie.\\
\end{itemize}

Lucrarea de disertație se tipărește față-verso.

Cuprinsul și listele tabelelor, figurilor și abrevierilor vor fi numerotate folosind cifre romane (prima pagină a cuprinsului: pagina v). Capitolele propriu-zise și bibliografia vor fi numerotate folosind cifre arabe (prima pagină a primului capitol: pagina 1).

Toate secțiunile (cuprinsul, listele tabelelor, figurilor și abrevierilor, capitolele și bibliografia) trebuie să înceapă pe pagină \textbf{nouă}, \textbf{dreaptă} (deci impară); dacă este necesar, se vor lăsa pagini goale și nenumerotate, după caz (exemplificat în prezentul ghid).

Se recomandă ca introducerea (capitolul 1) să cuprindă 2\textbf{--}4 pagini, concluziile (ultimul capitol) să cuprindă 2\textbf{--}4 pagini, iar restul capitolelor (capitolul 2, 3, ...) să fie proporționale. De asemenea, subcapitolele și subcapitolele secundare trebuie să fie și ele proporționale. Nu se justifică să existe capitole de 3 pagini față de capitole de 20 de pagini, la fel cum nu se justifică subcapitole mai scurte de 1 pagină sau subcapitole secundare mai scurte de jumătate de pagină.

Se recomandă ca lucrarea de disertație să aibă \textbf{între 60 și 80 de pagini}.



%\begin{figure}[h]
%    \centering
%    \includegraphics[width=0.5\textwidth]{figuri/punguta.jpg}
%    \caption{Exemplu de figură.}
%    \label{fig:exemplu}
%\end{figure}


%Figura \ref{fig:exemplu} reprezintă un exemplu.


